% ===========================================================================
% CUMCM 数学建模论文 LaTeX 模板
% 全国大学生数学建模竞赛 (CUMCM) 中文论文格式
%
% 格式规范:
%   - 摘要单独成页，不超过1页
%   - 三线表 (booktabs: toprule/midrule/bottomrule, 1.5磅, 无竖线)
%   - 图题在下，表题在上
%   - 图 600dpi, 黑白可区分 (线型+图例, 非颜色)
%   - 公式居中，右对齐编号
%   - 参考文献5-10篇，[J][M][D][EB/OL]
%   - 宋体正文 + Times New Roman 英文/数学
%   - 页边距 2.5cm
%   - 正文不超过20页
% ===========================================================================

\documentclass[12pt,a4paper]{article}

% ---------- 中文支持 ----------
\usepackage[UTF8,heading=true]{ctex}
\ctexset{
    section = {
        format = \heiti\zihao{4}\centering,
        beforeskip = 12pt,
        afterskip = 6pt,
    },
    subsection = {
        format = \heiti\zihao{-4},
        beforeskip = 8pt,
        afterskip = 4pt,
    },
}

% ---------- 页面设置: 2.5cm 页边距 ----------
\usepackage[top=2.5cm,bottom=2.5cm,left=2.5cm,right=2.5cm]{geometry}

% ---------- 三线表 (booktabs, 1.5磅) ----------
\usepackage{booktabs}
\renewcommand{\toprule}{\specialrule{1.5pt}{0pt}{0pt}}
\renewcommand{\midrule}{\specialrule{1.5pt}{0pt}{0pt}}
\renewcommand{\bottomrule}{\specialrule{1.5pt}{0pt}{0pt}}

% ---------- 数学与符号 ----------
\usepackage{amsmath,amssymb}

% ---------- 图形支持 ----------
\usepackage{graphicx}
\usepackage[font=small,labelfont=bf,labelsep=space]{caption}
% 图题在图片下方 (默认设置)
\captionsetup[figure]{position=below}
% 表题在表格上方
\captionsetup[table]{position=above}

% ---------- 超链接 ----------
\usepackage[colorlinks=true,linkcolor=black,citecolor=black,urlcolor=black]{hyperref}

% ---------- 参考文献 ----------
\usepackage[numbers,sort&compress]{natbib}

% ---------- 行距: 1.5倍 ----------
\usepackage{setspace}
\onehalfspacing

% ---------- 附录 ----------
\usepackage[titletoc]{appendix}

% ---------- 代码展示 ----------
\usepackage{listings}
\usepackage{xcolor}
\lstset{
    basicstyle=\ttfamily\small,
    breaklines=true,
    frame=single,
    numbers=left,
    numberstyle=\tiny,
}

% ===========================================================================
% 论文信息 (请在此填写)
% ===========================================================================
\newcommand{\cumcmtitle}{论文标题}
\newcommand{\cumcmkeywords}{关键词1；关键词2；关键词3；关键词4}

% ===========================================================================
% 正文开始
% ===========================================================================
\begin{document}

% ===================================================================
% 摘要页 (单独成页)
% ===================================================================
\begin{center}
    {\heiti\zihao{3} \cumcmtitle}

    \vspace{12pt}
    {\heiti\zihao{4} 摘\quad 要}
\end{center}

\vspace{6pt}

% --- 注意事项 ---
% 摘要要求:
%   1. 含5要素: 模型归类 / 建模思想 / 算法思想 / 模型特色 / 主要结果
%   2. 3-5个关键词
%   3. 不得包含公式、表格、图片
%   4. 不超过1页

本文针对\underline{\hspace{2cm}}问题，建立了\underline{\hspace{2cm}}模型。
首先，\underline{\hspace{4cm}}。
其次，\underline{\hspace{4cm}}。
最后，\underline{\hspace{4cm}}。

主要结果如下:
\begin{enumerate}
    \item 问题一: \underline{\hspace{6cm}}
    \item 问题二: \underline{\hspace{6cm}}
    \item 问题三: \underline{\hspace{6cm}}
\end{enumerate}

\vspace{12pt}
\noindent {\heiti 关键词:} \cumcmkeywords

\newpage

% ===================================================================
% 正文从此页开始 (不超过20页)
% ===================================================================

% ===================================================================
% 1. 问题重述
% ===================================================================
\section{问题重述}
% 将原问题用自己的语言重新描述，列出需解决的问题。

\subsection{问题背景}
% 问题背景介绍

\subsection{问题的提出}
% 具体问题表述

\subsection{需要解决的问题}
% 列出需要解决的所有问题

% ===================================================================
% 2. 问题分析
% ===================================================================
\section{问题分析}
% 对每个问题的分析思路、涉及的数学方法

\subsection{问题一的分析}
\subsection{问题二的分析}
\subsection{问题三的分析}

% ===================================================================
% 3. 模型假设与符号说明
% ===================================================================
\section{模型假设与符号说明}

\subsection{基本假设}
\begin{enumerate}
    \item 假设题目所给数据真实可靠。
    \item \underline{\hspace{6cm}}
    \item \underline{\hspace{6cm}}
\end{enumerate}

\subsection{符号说明}
% 使用三线表列出符号
\begin{table}[htbp]
    \centering
    \caption{符号说明}
    \label{tab:symbols}
    \begin{tabular}{ccc}
        \toprule
        符号 & 含义 & 单位 \\
        \midrule
        $x$ & \underline{\hspace{3cm}} & \underline{\hspace{2cm}} \\
        $y$ & \underline{\hspace{3cm}} & \underline{\hspace{2cm}} \\
        \bottomrule
    \end{tabular}
\end{table}

% ===================================================================
% 4. 模型建立与求解
% ===================================================================
\section{模型建立与求解}

\subsection{模型一: \underline{\hspace{3cm}}}
% 模型建立
% 公式 (居中, 后文引用才编号)
% 求解方法
% 结果
% 结果展示用三线表 (表题在上)

\begin{table}[htbp]
    \centering
    \caption{模型一结果}
    \label{tab:model1}
    \begin{tabular}{ccc}
        \toprule
        指标 & 数值 & 备注 \\
        \midrule
        \underline{\hspace{2cm}} & \underline{\hspace{2cm}} & \underline{\hspace{2cm}} \\
        \bottomrule
    \end{tabular}
\end{table}

% 示例公式 (后文引用才编号)
\begin{equation}
    y = f(x) = ax + b
    \label{eq:linear}
\end{equation}

\subsection{模型二: \underline{\hspace{3cm}}}
% 模型建立、求解、结果

% ===================================================================
% 5. 模型评价与改进
% ===================================================================
\section{模型评价与改进}

\subsection{模型的优点}
\subsection{模型的缺点}
\subsection{模型的改进方向}

% ===================================================================
% 6. 模型推广
% ===================================================================
\section{模型推广}

% ===================================================================
% 参考文献 (5-10篇)
% 格式: [J]期刊文章 [M]专著 [D]学位论文 [EB/OL]电子文献
% 使用英文标点+空格
% ===================================================================
\begin{thebibliography}{99}

\bibitem{ref1} 作者1, 作者2. 文章标题[J]. 期刊名, 年份, 卷(期): 页码.

\bibitem{ref2} 作者. 书名[M]. 出版地: 出版社, 年份.

\bibitem{ref3} 作者. 论文标题[D]. 学校, 年份.

\bibitem{ref4} 作者. 文献标题[EB/OL]. URL, 访问日期.

\bibitem{ref5} 作者1, 作者2. 文章标题[J]. 期刊名, 年份, 卷(期): 页码.

\end{thebibliography}

% ===================================================================
% 附录 (单独成页, 含全部可运行代码)
% ===================================================================
\newpage
\appendix
\section{附录A: 程序代码}

% ------ 示例代码 ------
% \begin{lstlisting}[caption={主程序},label=lst:main]
% import numpy as np
% import pandas as pd
% # ... your code here
% \end{lstlisting}

\end{document}
